\documentclass[11pt]{article}
\usepackage{fontspec,amsmath,amssymb,amsthm,hyperref}
\usepackage[margin=1in]{geometry}
\setlength{\emergencystretch}{3em}
\title{The complete first-order KLMC theorem for the actual sampler}
\author{LangevinLean mathematical exposition}
\date{}
\begin{document}\maketitle
\noindent Source module: \texttt{KineticLangevin/CanonicalKLMC.lean}. Compile this proof document with XeLaTeX or LuaLaTeX. The accompanying HTML pairs the mathematical argument with the exact Lean statements and records the scope of the formalization.


This final theorem concerns the implemented KLMC chain on the explicit independent Gaussian-velocity and canonical Brownian source. It proves the first-order Wasserstein bound of Dalalyan and Riou-Durand with the true initial Wasserstein distance. All continuous reference laws, common-noise identities, local-error estimates, output-law identities and transport-infimum steps are discharged inside the proof chain.


\section{Prove Gibbs stationarity of the actual global coupling reference}


The reference begins with standard Gaussian velocity independent of the target Gibbs position. Its full Brownian forcing and its entire filtered innovation stream are actual constructed objects on the initial-coupling extension. Their exact initial/noise product law is already proved. The kinetic stationarity transfer theorem therefore applies directly to the actual integral solution. For a time \(t\) inside a larger horizon, Lean restricts the actual integral equation to \([0,t]\), supplies the continuous input and its zero starting value, and applies the proved splitting-limit stationarity result. This yields the full phase law \(\mathcal N(0,I_d)\otimes\pi\) at every reference time, without a remaining stationarity or Gaussian-marginal hypothesis.



\noindent Lean declarations: \texttt{globalReferenceFlow\_stationary}.

\section{Conclude the source-paper KLMC Wasserstein estimate}


Let \(f\) be differentiable, \(m\)-strongly convex with \(m>0\), and have an \(M\)-Lipschitz gradient. If \(\gamma^2\ge m+M\), \(\gamma>0\), and \(0<h\le m/(4\gamma M)\), the actual KLMC position law satisfies \[W_2(\nu_n,\pi)\le\sqrt2\left(1-\frac{3mh}{4\gamma}\right)^nW_2(\mu,\pi)+\frac{Mh\sqrt{2d}}m.\] The initial position distribution may be arbitrary with finite initial Wasserstein distance; its velocity is the explicit independent standard Gaussian. Lean first derives the finite Gibbs second moment from strong convexity. It then supplies the just-proved global-reference stationarity for every initial transport coupling to the complete canonical coupling/infimum theorem. The conclusion has no conditional local-error, reference-process, Gaussian-law, stationary-law or optimal-coupling premise. The smoothness assumptions are expressed by the actual gradient and strong convexity, so the theorem also covers the paper's Hessian-bounded setting through the existing deterministic geometry results.



\noindent Lean declarations: \texttt{klmc\_wasserstein\_bound}.

\end{document}
